% ============================================================
%  70asunflower Blog Template
%  通用技术博客 LaTeX 模板
%  Compile with: XeLaTeX
% ============================================================
\documentclass[fontset=fandol,11pt,a4paper]{ctexart}

% 消除 Fandol 字体 Bold/Italic 警告
\usepackage{xeCJK}
\xeCJKsetup{AutoFakeBold=true, AutoFakeSlant=true}

% ---------- 基础包 ----------
\usepackage{geometry}
\usepackage{xcolor}
\usepackage{amsmath,amssymb,amsfonts,amsthm}
\usepackage{booktabs}
\usepackage{array}
\usepackage{enumitem}
\usepackage{tikz}
\usepackage{tikzpagenodes}
\usepackage{fancyhdr}
\usepackage{titlesec}
\usepackage{titletoc}
\usepackage{etoolbox}
\usepackage{hyperref}
\usepackage{graphicx}
\usepackage{eso-pic}
\usepackage{listings}

% ---------- 页面设置 ----------
\geometry{
  left=2.4cm, right=2.4cm,
  top=3.2cm, bottom=2.5cm,
  headheight=25pt
}

% ---------- 配色 ----------
\definecolor{Primary}{RGB}{42, 82, 120}
\definecolor{Secondary}{RGB}{232, 119, 34}
\definecolor{Accent}{RGB}{142, 68, 173}
\definecolor{DarkText}{RGB}{44, 62, 80}
\definecolor{LightBg}{RGB}{250, 248, 245}
\definecolor{CodeBg}{RGB}{245, 245, 240}
\definecolor{TipBg}{RGB}{232, 245, 233}
\definecolor{WarnBg}{RGB}{255, 243, 224}
\definecolor{FormulaBg}{RGB}{240, 244, 248}
\definecolor{SoftBorder}{RGB}{220, 215, 210}

% ---------- listings 代码样式 ----------
\lstset{
  basicstyle=\ttfamily\small\color{DarkText},
  backgroundcolor=\color{CodeBg},
  frame=none,
  breaklines=true,
  showstringspaces=false,
  tabsize=2,
  keywordstyle=\color{Primary}\bfseries,
  commentstyle=\color{gray!60}\itshape,
  stringstyle=\color{Secondary},
  numbers=none,
  xleftmargin=0pt,
  xrightmargin=0pt,
  aboveskip=0pt,
  belowskip=0pt
}

% ---------- hyperref ----------
\hypersetup{
  colorlinks=true,
  linkcolor=Primary,
  citecolor=Accent,
  urlcolor=Secondary,
  pdftitle={70asunflower Blog},
  pdfauthor={70asunflower}
}

% ---------- 水印：向日葵（纯TikZ绘制）----------
\newcommand{\sunflowerwatermark}[2][0.15]{
  \begin{tikzpicture}[scale=#2]
    \foreach \i in {0,22.5,...,337.5} {
      \fill[Secondary!#1!yellow, rotate=\i] (0,0) ellipse (0.6 and 0.25);
    }
    \fill[Primary!#1!brown] (0,0) circle (0.3);
    \fill[Secondary!#1!orange] (0,0) circle (0.12);
  \end{tikzpicture}
}

% 每页背景水印
\AddToShipoutPictureBG{
  \begin{tikzpicture}[remember picture, overlay]
    \node[opacity=0.03, rotate=25] at (current page.center) {
      \fontsize{55}{66}\selectfont\color{Secondary}\textbf{70asunflower}
    };
    \node[opacity=0.06] at ([xshift=-1.5cm, yshift=2cm]current page.south east) {
      \sunflowerwatermark[30]{0.8}
    };
  \end{tikzpicture}
}

% ---------- 页眉页脚 ----------
\pagestyle{fancy}
\fancyhf{}
\fancyhead[L]{\small\color{Primary}$\star$ 70asunflower}
\fancyhead[R]{\small\color{Primary}\thesection\ \leftmark}
\fancyfoot[C]{\small\color{gray!50}--- \thepage\ ---}
\renewcommand{\headrulewidth}{0.4pt}
\renewcommand{\headrule}{\hbox to\headwidth{\color{SoftBorder}\leaders\hrule height \headrulewidth\hfill}}
\renewcommand{\footrulewidth}{0pt}
\renewcommand{\sectionmark}[1]{\markboth{#1}{}}

% ---------- 章节标题 ----------
\titleformat{\section}
  {\normalfont\LARGE\bfseries\color{Primary}}
  {\color{Secondary}\thesection}{0.8em}{}
  [\vspace{-0.5em}\textcolor{SoftBorder}{\rule{\textwidth}{1pt}}]

\titleformat{\subsection}
  {\normalfont\large\bfseries\color{Primary!80}}
  {\color{Secondary}\thesubsection}{0.8em}{}

\titleformat{\subsubsection}
  {\normalfont\normalsize\bfseries\color{Primary!60}}
  {\color{Secondary!80}\thesubsubsection}{0.8em}{}

% ---------- 目录 ----------
\titlecontents{section}[0pt]
  {\bfseries\color{Primary}\addvspace{6pt}}
  {\contentslabel{1.8em}}
  {}
  {\color{Secondary}\titlerule*[6pt]{.}\contentspage}
\titlecontents{subsection}[1.8em]
  {\color{Primary!80}}
  {\contentslabel{2.2em}}
  {}
  {\color{SoftBorder}\titlerule*[6pt]{.}\contentspage}

% ---------- tcolorbox ----------
\usepackage[most,breakable,skins]{tcolorbox}

% 封面卡片
\newtcolorbox{covercard}[1][]{
  enhanced, colback=LightBg, colframe=Primary, arc=4mm, boxrule=0pt,
  leftrule=4pt, left=20pt, right=20pt, top=18pt, bottom=18pt,
  shadow={2mm}{-2mm}{0mm}{black!20}, #1
}

% 摘要框
\newtcolorbox{abstractcard}[1][]{
  enhanced, colback=LightBg, colframe=SoftBorder, arc=2mm, boxrule=0pt,
  leftrule=4pt, left=14pt, right=14pt, top=12pt, bottom=12pt,
  borderline west={4pt}{0pt}{Secondary},
  fonttitle=\bfseries\color{Primary}, coltitle=Primary,
  attach boxed title to top left={yshift=-2mm, xshift=8mm},
  boxed title style={colback=LightBg, arc=1mm, boxrule=0pt, left=8pt, right=8pt},
  title={$\diamond$ 摘要}, #1
}

% 公式框
\newtcolorbox{mathbox}[1][]{
  enhanced, colback=FormulaBg, colframe=FormulaBg, arc=3mm, boxrule=0pt,
  left=16pt, right=16pt, top=14pt, bottom=14pt,
  shadow={1mm}{-1mm}{0mm}{black!10}, #1
}

% 代码框
\newtcolorbox{codebox}[1][Python]{
  enhanced, colback=CodeBg, colframe=SoftBorder, arc=2mm, boxrule=0.5pt,
  left=12pt, right=12pt, top=10pt, bottom=10pt,
  fonttitle=\bfseries\small\color{Primary}, coltitle=Primary,
  attach boxed title to top left={yshift=-3mm, xshift=6mm},
  boxed title style={colback=CodeBg, arc=1mm, boxrule=0.5pt, left=8pt, right=8pt, top=1pt, bottom=1pt},
  title={$\triangleright$ #1},
}

% 提示框
\newtcolorbox{tipbox}[1][]{
  enhanced, colback=TipBg, colframe=TipBg, arc=2mm, boxrule=0pt,
  left=12pt, right=12pt, top=10pt, bottom=10pt,
  borderline west={3pt}{0pt}{green!60!black}, fontupper=\small, #1
}

% 警告框
\newtcolorbox{warnbox}[1][]{
  enhanced, colback=WarnBg, colframe=WarnBg, arc=2mm, boxrule=0pt,
  left=12pt, right=12pt, top=10pt, bottom=10pt,
  borderline west={3pt}{0pt}{Secondary}, fontupper=\small, #1
}

% 引用框
\newtcolorbox{quotebox}[1][]{
  enhanced, colback=LightBg, colframe=SoftBorder, arc=1mm, boxrule=0pt,
  leftrule=3pt, left=14pt, right=14pt, top=8pt, bottom=8pt,
  fontupper=\itshape\color{DarkText!80}, #1
}

% ---------- 列表 ----------
\setlist[itemize]{leftmargin=*,label=\textcolor{Secondary}{\textbullet},topsep=4pt,itemsep=3pt}
\setlist[enumerate]{leftmargin=*,topsep=4pt,itemsep=3pt}

% ---------- 自定义命令 ----------
\newcommand{\keyword}[1]{\textbf{\color{Primary}#1}}
\newcommand{\concept}[1]{\textbf{\color{Accent}#1}}
\newcommand{\highlight}[1]{\colorbox{Secondary!15}{#1}}
\newcommand{\codeinline}[1]{\texttt{\color{Primary!80}#1}}

% ---------- 文档开始 ----------
\begin{document}

% ==================== 封面（AI修改位置）====================
\thispagestyle{empty}

\begin{tikzpicture}[remember picture, overlay]
  \fill[Primary] (current page.north west) rectangle ([yshift=-0.6cm]current page.north east);
  \fill[Primary!80] (current page.south west) rectangle ([yshift=0.3cm]current page.south east);
\end{tikzpicture}

\vspace*{1cm}

\begin{center}
  {\sunflowerwatermark[100]{2.5}}

  \vspace{0.4cm}

  % AI修改：日期、作者、标签
  {\small\color{gray!70}$\bullet$ 2026年6月 \quad $\bullet$ 70asunflower \quad $\bullet$ 模板, 教程}

  \vspace{0.6cm}

  \begin{covercard}
    % AI修改：标题
    {\Huge\bfseries\color{Primary} 70asunflower 博客模板}

    \vspace{0.3cm}

    % AI修改：副标题
    {\Large\color{Primary!80} Markdown 转 LaTeX 格式规范手册}

    \vspace{0.4cm}

    {\normalsize\color{gray!70}$\star$ 70asunflower \quad 技术博客}
  \end{covercard}

  \vspace{0.6cm}

  % AI修改：摘要
  \begin{abstractcard}
    本文档是模板的自解释手册。AI 读取此模板后，应能根据以下规则将 Markdown 内容自动转换为该 LaTeX 格式。
  \end{abstractcard}
\end{center}

\vspace{0.8cm}

% ==================== 目录（罗马数字页码，避免与正文冲突）====================
\pagenumbering{Roman}
\tableofcontents

% ==================== 正文：AI转换规则手册 ====================
\newpage
\pagenumbering{arabic}
\setcounter{page}{2}

\section{封面与元信息}
\label{sec:meta}

\subsection{修改位置说明}
\label{subsec:meta-edit}

AI 生成封面时，需替换以下四处内容：

\begin{center}
\begin{tabular}{>{\color{Primary}}l >{\color{DarkText}}l >{\color{gray!70}}l}
\toprule
\textbf{位置} & \textbf{LaTeX 代码} & \textbf{对应 Markdown} \\
\midrule
标题 & \verb|{\Huge\bfseries\color{Primary} 标题}| & \verb|# 标题| \\
副标题 & \verb|{\Large\color{Primary!80} 副标题}| & 首段描述或无 \\
日期 & \verb|$\bullet$ 2026年6月| & YAML frontmatter \verb|date| \\
作者 & \verb|$\bullet$ 70asunflower| & YAML frontmatter \verb|author| \\
标签 & \verb|$\bullet$ 标签1, 标签2| & YAML frontmatter \verb|tags| \\
摘要 & \verb|\begin{abstractcard} ... \end{abstractcard}| & \verb|<!-- abstract -->| \\
\bottomrule
\end{tabular}
\end{center}

\begin{warnbox}
\textbf{$\triangle$ 注意：} 封面页使用 \verb|\thispagestyle{empty}|，页码从目录后的正文开始计数。
\end{warnbox}

\section{章节与目录}
\label{sec:section}

\subsection{章节层级映射}
\label{subsec:section-map}

\begin{center}
\begin{tabular}{>{\color{Primary}}l >{\color{DarkText}}l}
\toprule
\textbf{Markdown} & \textbf{LaTeX} \\
\midrule
\verb|# 一级标题| & \verb|\section{一级标题}| \\
\verb|## 二级标题| & \verb|\subsection{二级标题}| \\
\verb|### 三级标题| & \verb|\subsubsection{三级标题}| \\
\bottomrule
\end{tabular}
\end{center}

\begin{tipbox}
\textbf{$\rightarrow$ 提示：} 目录由 \verb|\tableofcontents| 自动生成，无需手动维护。章节标题会自动显示在页眉右侧。
\end{tipbox}

\section{公式排版}
\label{sec:formula}

\subsection{行内公式}
\label{subsec:inline-formula}

Markdown 中的 \verb|$...$| 直接保留为 LaTeX 行内公式：
\[
\text{Markdown: } E=mc^2 \quad \Rightarrow \quad \text{LaTeX: } \verb|$E=mc^2$|
\]

\subsection{独立公式（放入公式框）}
\label{subsec:display-formula}

Markdown 中的 \verb|$$...$$| 或代码块标记为 \verb|math| 时，应放入 \verb|mathbox| 环境：

\begin{mathbox}
\[
\int_{-\infty}^{+\infty} e^{-x^2} dx = \sqrt{\pi}
\]
\end{mathbox}

对应 LaTeX 代码：
\begin{codebox}[LaTeX]
\begin{lstlisting}[language=TeX]
\begin{mathbox}
\[
\int_{-\infty}^{+\infty} e^{-x^2} dx = \sqrt{\pi}
\]
\end{mathbox}
\end{lstlisting}
\end{codebox}

\section{代码块}
\label{sec:code}

\subsection{基本规则}
\label{subsec:code-rule}

Markdown 中的 \verb| ```python | 代码块，应转换为 \verb|codebox| + \verb|lstlisting| 环境：

\begin{codebox}[Python]
\begin{lstlisting}[language=Python]
def hello_world(name: str) -> str:
    """返回问候语"""
    return f"Hello, {name}!"
\end{lstlisting}
\end{codebox}

\subsection{语言标签映射}
\label{subsec:code-lang}

\begin{center}
\begin{tabular}{>{\color{Primary}}l >{\color{DarkText}}l}
\toprule
\textbf{Markdown 标记} & \textbf{codebox 标题参数} \\
\midrule
\verb|python| & \verb|[Python]| \\
\verb|bash| / \verb|shell| & \verb|[Bash]| \\
\verb|javascript| / \verb|js| & \verb|[JavaScript]| \\
\verb|cpp| / \verb|c++| & \verb|[C++]| \\
\verb|tex| / \verb|latex| & \verb|[LaTeX]| \\
其他 & \verb|[Code]| \\
\bottomrule
\end{tabular}
\end{center}

\begin{warnbox}
\textbf{$\triangle$ 注意：} 切勿在 \verb|codebox| 内部使用 \verb|verbatim| 环境，必须使用 \verb|lstlisting|，否则编译报错。
\end{warnbox}

\section{特殊框体}
\label{sec:boxes}

\subsection{提示框（Tip）}
\label{subsec:tip}

Markdown 中的引用块若包含 \textbf{提示} 关键词，转换为 \verb|tipbox|：

\begin{tipbox}
\textbf{$\rightarrow$ 提示：} 这是提示框的内容。左侧有绿色竖线标识。
\end{tipbox}

LaTeX 代码：
\begin{codebox}[LaTeX]
\begin{lstlisting}[language=TeX]
\begin{tipbox}
\textbf{$\rightarrow$ 提示：} 内容...
\end{tipbox}
\end{lstlisting}
\end{codebox}

\subsection{警告框（Warning）}
\label{subsec:warn}

Markdown 中的引用块若包含 \textbf{警告/注意} 关键词，转换为 \verb|warnbox|：

\begin{warnbox}
\textbf{$\triangle$ 注意：} 这是警告框的内容。左侧有橙色竖线标识。
\end{warnbox}

\subsection{引用框（Quote）}
\label{subsec:quote}

普通引用块转换为 \verb|quotebox|：

\begin{quotebox}
这是引用框的内容，默认斜体显示，适合放置重要结论或假设。
\end{quotebox}

\section{列表与表格}
\label{sec:list-table}

\subsection{无序列表}
\label{subsec:ul}

Markdown：
\begin{codebox}[Markdown]
\begin{lstlisting}
- 项目一
- 项目二
- 项目三
\end{lstlisting}
\end{codebox}

LaTeX 转换结果：
\begin{itemize}
  \item 项目一
  \item 项目二
  \item 项目三
\end{itemize}

\subsection{有序列表}
\label{subsec:ol}

Markdown：
\begin{codebox}[Markdown]
\begin{lstlisting}
1. 步骤一
2. 步骤二
3. 步骤三
\end{lstlisting}
\end{codebox}

LaTeX 转换结果：
\begin{enumerate}
  \item 步骤一
  \item 步骤二
  \item 步骤三
\end{enumerate}

\subsection{表格}
\label{subsec:table}

Markdown 表格转换为 \verb|booktabs| 风格的 \verb|tabular|：

\begin{center}
\begin{tabular}{lll}
\toprule
\textbf{姓名} & \textbf{年龄} & \textbf{职业} \\
\midrule
Alice & 24 & 工程师 \\
Bob & 30 & 研究员 \\
\bottomrule
\end{tabular}
\end{center}

\section{文本高亮}
\label{sec:highlight}

\subsection{关键词与概念}
\label{subsec:keyword}

\begin{center}
\begin{tabular}{>{\color{Primary}}l >{\color{DarkText}}l >{\color{gray!70}}l}
\toprule
\textbf{Markdown} & \textbf{LaTeX 命令} & \textbf{效果} \\
\midrule
\verb|**关键词**| & \verb|\keyword{关键词}| & \keyword{蓝色粗体} \\
\verb|*概念*| & \verb|\concept{概念}| & \concept{紫色粗体} \\
\verb|==高亮==| & \verb|\highlight{高亮}| & \highlight{橙色背景} \\
\verb|``代码``| & \verb|\codeinline{代码}| & \codeinline{等宽字体} \\
\bottomrule
\end{tabular}
\end{center}

\section{结语与页脚}
\label{sec:footer}

\subsection{结语写法}
\label{subsec:conclusion}

文档结尾使用无编号章节 \verb|\section*{结语}|，并放入 \verb|quotebox|：

\begin{quotebox}
本文档展示了 70asunflower 博客模板的全部转换规则。AI 读取此模板后，应能自动化完成 Markdown 到 LaTeX 的格式转换。
\end{quotebox}

\subsection{页脚水印}
\label{subsec:watermark}

每页背景自动包含 \textbf{70asunflower} 文字水印和右下角向日葵图标。结语页可额外放置大图标：

\begin{center}
{\sunflowerwatermark[80]{1.5}}
\end{center}

\vspace{1cm}

\begin{center}
\textcolor{gray!60}{\small$\star$ 70asunflower \quad $\heartsuit$ 感谢阅读，欢迎交流 \quad $\star$ 全网同名}
\end{center}

\end{document}
